\section{A workload-aware algorithm-selection policy}\label{sec:policy}
The theoretical problem is to choose a complete pipeline, not a favored numerical library in isolation. For a specified tolerance $\varepsilon$, memory budget $B$, difference order $d$, number of candidates $M$, observed-series count $R$ and criterion set $\mathcal C$, we recommend
\begin{equation}\label{eq:frontier}
\mathcal A^\star\in
\arg\min_{\mathcal A\in\mathfrak A_{\rm valid}}
T_\mathcal A(n,d,M,R,\mathcal C;\varepsilon,B),
\end{equation}
where the constraint \eqref{eq:valid} is evaluated before comparing cost. Neither Theorem~\ref{thm:svdwork} nor Theorem~\ref{thm:bandwork} supplies measured wall-clock constants, so \eqref{eq:frontier} is a conditional decision principle, not an empirically calibrated universal dispatcher.

\begin{table*}[t]
\caption{Algorithm choices suggested by the derived cost components. Each preference is conditional on accurate outputs, an admissible $S$ inversion and the complete criterion being available.}
\label{tab:policy}\centering\small
\setlength{\tabcolsep}{5pt}
\begin{tabularx}{\textwidth}{@{}p{.24\textwidth}Xp{.26\textwidth}@{}}
\toprule
Workload & Initial candidate and rationale & Important qualification\\
\midrule
One curve, known $\lambda$, small fixed $d$ & Banded Cholesky: $O(nd^2)$ factor, $O(nd)$ solve & Large $\lambda$ can amplify finite-precision error.\\
One curve, prescribed $S$ & Banded fit with a certified $S\mapsto\lambda$ mapper & Do not omit mapping cost. For $d=1$ the spectrum is analytic.\\
Many $S$ values, scalar GCV or likelihood, moderate $n$ & Retained SVD: costly setup, $O(r)$ candidate scores & Dense basis must fit memory; finite constants can reverse comparisons.\\
Many $S$ values, exact ordinary CV, small $d$ and large $n$ & Banded Cholesky plus selected inverse & Spectral CV also has $O(nr)$ leverage/residual reconstruction.\\
Many series sharing $(n,d,\mathcal S)$ & Map $S$ once; reuse spectrum or factors and leverage diagnostics & Each series still needs its own projections or solves.\\
Near $S=1$, large $d$, difficult conditioning & Validated shift-free spectral or independently verified structured alternative & Prefer analytic endpoint at exact $S=1$; never assume a particular algorithm always passes.\\
\bottomrule
\end{tabularx}
\end{table*}

\subsection{Complete candidate-grid procedure}
The following abstract algorithm specifies a reproducible computation rather than claiming that an optimal dispatcher is already implemented in the library.
\begin{enumerate}[leftmargin=*]
\item Fix $(n,d)$, the reference-mean policy, the criterion set $\mathcal C$, candidate set $\mathcal S$, numerical tolerances and memory limit. Declare whether any mean estimate is conditional or foldwise-refitted.
\item Determine the exact admissible endpoint conventions. Map each distinct interior $s_j$ to a penalty once, checking its inverse-map residual. Do not independently space a comparison grid in $\log\lambda$ and call it the same statistical workload.
\item Evaluate complete pipeline feasibility. A spectral method must budget its full basis; a structured method must budget all selected-inverse and likelihood diagnostics, not merely the trend fit.
\item Reuse calculations at the correct level: the $S$ mapping is shared across all series of the same length and order; a banded factor or leverage vector is shared across series for the same penalty; spectral projections depend on the observed series and fixed mean.
\item At each candidate, evaluate all requested scores and numerical checks. Record invalid or undefined scores and failed fits rather than assigning them artificial best values.
\item Return a minimizer \emph{over the explicitly evaluated finite grid} for each requested criterion. The result is not automatically a certified global minimizer of the unrestricted continuous criterion.
\end{enumerate}

\subsection{What can be concluded without empirical timing}
The theory proves that structured factorization has favorable asymptotic storage and single-fit arithmetic for fixed small $d$. It also proves that once an SVD has been paid for, many scalar smoothing scores are cheap and share the same spectral basis and smoothness map. It does \emph{not} establish that SVD always beats banded Cholesky on CV grids, that a positive crossover exists in every regime, or that a faster measured fit implies faster complete criterion selection. Those are contingent questions involving operation constants and numerical admissibility.

Finally, changing $d$ is a model-design decision, not merely a computational tuning parameter. Different difference orders impose different nullspaces and different smoothing trajectories. A preference inferred from computational work must not be misreported as a guarantee of better latent-trend recovery or forecasting accuracy.
